\subsection{Berechnung einer OPV-Schaltung mit drei Eingängen}
% Alt: Aufgabe 16
\Aufgabe{
    \label{Aufg16}
    Gegeben ist die unten abgebildete Schaltung. Bestimmen Sie die Ausgangsspannung \( U_\mathrm{A} \). \\
    Die Widerstandswerte lauten: \\
    $
        R_\mathrm{1} = \mathrm{10\,k\Omega},\
        R_\mathrm{2} = \mathrm{20\,k\Omega},\
        R_\mathrm{3} = \mathrm{30\,k\Omega},\
        R_\mathrm{4} = \mathrm{10\,k\Omega},\
        R_\mathrm{5} = \mathrm{20\,k\Omega},\
        R_\mathrm{6} = \mathrm{40\,k\Omega}
    $
    \begin{figure}[H]
        \centering
        \resizebox{0.95\textwidth}{!}{\input{Tikz/src/FigOPVzweistufig.tex}\alttext{\alttext{Das Bild zeigt einen elektronischen Schaltkreis mit sechs Widerständen und zwei Operationsverstärkern. Links sind drei Spannungsquellen mit den Bezeichnungen \(U_{\mathrm{E1}}\), \(U_{\mathrm{E2}}\) und \(U_{\mathrm{E3}}\), die von oben nach unten angeordnet sind. Die Widerstände sind beschriftet als \(R_1\), \(R_2\), \(R_3\), \(R_4\), \(R_5\) und \(R_6\).
                    \(R_1\) ist oben links parallel zu \(R_2\) darunter geschaltet. Darunter befindet sich \(R_3\), der mit einem Knotenpunkt verbunden ist. Von diesem Knoten führt eine Leitung zu dem ersten Operationsverstärker N1. Zwischen dem Eingang des Operationsverstärkers N1 und seinem Ausgang liegt der Widerstand \(R_4\) in einer Rückkopplungsschleife. Der Ausgang von N1 ist über den Widerstand \(R_5\) mit dem Eingang des zweiten Operationsverstärkers N2 verbunden.
                    Zwischen dem Eingang und Ausgang von N2 befindet sich der Widerstand \(R_6\) in einer Rückkopplungsschleife. Der Ausgang des zweiten Operationsverstärkers ist als Spannung \(U_\mathrm{A}\) bezeichnet. Alle Verbindungen sind durch Leitungen dargestellt, die die Bauteile miteinander verbinden und auf Masse bezogen sind.}}
        }

        \label{fig:FigOPV3Eingaenge}
    \end{figure}
}

\Loesung{
    Ausgangsspannung nach 1. OPV mit folgender Formel berechnen ($U_{\mathrm{A1}}$):
    \begin{align*}
        U_\mathrm{A1} & = - (U_\mathrm{E1} \cdot \frac{R_4}{R_1} + U_\mathrm{E2} \cdot \frac{R_4}{R_2} + U_\mathrm{E3} \cdot \frac{R_4}{R_3}) \\
        U_\mathrm{A1} & = - (U_\mathrm{E1} \cdot \frac{10}{10} + U_\mathrm{E2} \cdot \frac{10}{20} + U_\mathrm{E3} \cdot \frac{10}{30})       \\
        U_\mathrm{A1} & = - (U_\mathrm{E1} + \frac{1}{2} \cdot U_\mathrm{E2} + \frac{1}{3} \cdot U_\mathrm{E3})                               \\
    \end{align*}
    Daraus egibt sich für die gesamte Ausgangsspannung ($U_\mathrm{A}$):
    \begin{align*}
        U_\mathrm{A} & = - \frac{R_6}{R_5} \cdot U_\mathrm{\mathrm{A1}}                                                                \\
        U_\mathrm{A} & = - \frac{40}{20} \cdot (- (U_\mathrm{E1} + \frac{1}{2} \cdot U_\mathrm{E2} + \frac{1}{3} \cdot U_\mathrm{E3})) \\
        U_\mathrm{A} & = 2 \cdot U_\mathrm{E1} + U_\mathrm{E2} + \frac{2}{3} \cdot U_\mathrm{E3}
    \end{align*}


    Siehe: Abschnitt \ref{fig: Verschaltung Verstaerker} und Tabelle \ref{tab:Grundschaltungen}
}